%%%%%%%% ICML 2026 EXAMPLE LATEX SUBMISSION FILE %%%%%%%%%%%%%%%%%

\documentclass{article}

% Recommended, but optional, packages for figures and better typesetting:
\usepackage{microtype}
\usepackage{graphicx}
\usepackage{subcaption}
\usepackage{booktabs} % for professional tables
\usepackage{multirow}
\usepackage{xcolor}

% hyperref makes hyperlinks in the resulting PDF.
% If your build breaks (sometimes temporarily if a hyperlink spans a page)
% please comment out the following usepackage line and replace
% \usepackage{icml2026} with \usepackage[nohyperref]{icml2026} above.
\usepackage{hyperref}


% Attempt to make hyperref and algorithmic work together better:
\newcommand{\theHalgorithm}{\arabic{algorithm}}

% Use the following line for the initial blind version submitted for review:
\usepackage{icml2026}

% For preprint, use
% \usepackage[preprint]{icml2026}

% If accepted, instead use the following line for the camera-ready submission:
% \usepackage[accepted]{icml2026}

\usepackage{amsmath}
\usepackage{amssymb}
\usepackage{mathtools}
\usepackage{amsthm}


% if you use cleveref..
\usepackage[capitalize,noabbrev]{cleveref}

%%%%%%%%%%%%%%%%%%%%%%%%%%%%%%%%
% THEOREMS
%%%%%%%%%%%%%%%%%%%%%%%%%%%%%%%%
\theoremstyle{plain}
\newtheorem{theorem}{Theorem}[section]
\newtheorem{proposition}[theorem]{Proposition}
\newtheorem{lemma}[theorem]{Lemma}
\newtheorem{corollary}[theorem]{Corollary}
\theoremstyle{definition}
\newtheorem{definition}[theorem]{Definition}
\newtheorem{assumption}[theorem]{Assumption}
\theoremstyle{remark}
\newtheorem{remark}[theorem]{Remark}

\newcommand{\agent}[1]{}
\newcommand{\adam}[1]{}

% Todonotes is useful during development; simply uncomment the next line
%    and comment out the line below the next line to turn off comments
\usepackage[disable,textsize=tiny]{todonotes}

% The \icmltitle you define below is probably too long as a header.
% Therefore, a short form for the running title is supplied here:
\icmltitlerunning{Agentic Testing: Validating Informal Properties of Repository-Level Code}

\begin{document}

\twocolumn[
  % \icmltitle{Agentic Testing: Validating Informal Properties of Repository-Level Code}
  \icmltitle{Agentic Testing: Validating Repositories Against\\Natural Language Specifications}

  % It is OKAY to include author information, even for blind submissions: the
  % style file will automatically remove it for you unless you've provided
  % the [accepted] option to the icml2026 package.

  % List of affiliations: The first argument should be a (short) identifier you
  % will use later to specify author affiliations Academic affiliations
  % should list Department, University, City, Region, Country Industry
  % affiliations should list Company, City, Region, Country

  % You can specify symbols, otherwise they are numbered in order. Ideally, you
  % should not use this facility. Affiliations will be numbered in order of
  % appearance and this is the preferred way.
  \icmlsetsymbol{equal}{*}

  \begin{icmlauthorlist}
    \icmlauthor{Firstname1 Lastname1}{equal,yyy}
    \icmlauthor{Firstname2 Lastname2}{equal,yyy,comp}
    \icmlauthor{Firstname3 Lastname3}{comp}
    \icmlauthor{Firstname4 Lastname4}{sch}
    \icmlauthor{Firstname5 Lastname5}{yyy}
    \icmlauthor{Firstname6 Lastname6}{sch,yyy,comp}
    \icmlauthor{Firstname7 Lastname7}{comp}
    %\icmlauthor{}{sch}
    \icmlauthor{Firstname8 Lastname8}{sch}
    \icmlauthor{Firstname8 Lastname8}{yyy,comp}
    %\icmlauthor{}{sch}
    %\icmlauthor{}{sch}
  \end{icmlauthorlist}

  \icmlaffiliation{yyy}{Department of XXX, University of YYY, Location, Country}
  \icmlaffiliation{comp}{Company Name, Location, Country}
  \icmlaffiliation{sch}{School of ZZZ, Institute of WWW, Location, Country}

  \icmlcorrespondingauthor{Firstname1 Lastname1}{first1.last1@xxx.edu}
  \icmlcorrespondingauthor{Firstname2 Lastname2}{first2.last2@www.uk}

  % You may provide any keywords that you find helpful for describing your
  % paper; these are used to populate the "keywords" metadata in the PDF but
  % will not be shown in the document
  \icmlkeywords{Machine Learning, ICML}

  \vskip 0.3in
]

% this must go after the closing bracket ] following \twocolumn[ ...

% This command actually creates the footnote in the first column listing the
% affiliations and the copyright notice. The command takes one argument, which
% is text to display at the start of the footnote. The \icmlEqualContribution
% command is standard text for equal contribution. Remove it (just {}) if you
% do not need this facility.

% Use ONE of the following lines. DO NOT remove the command.
% If you have no special notice, KEEP empty braces:
\printAffiliationsAndNotice{}  % no special notice (required even if empty)
% Or, if applicable, use the standard equal contribution text:
% \printAffiliationsAndNotice{\icmlEqualContribution}

\begin{abstract}
Software correctness remains difficult to validate at the repository level, especially when intent is specified only in natural language.
At the same time, highly capable LLM agents increasingly modify codebases from high-level instructions.
Given the high coding capabilities of today's LLM agents, we propose Agentic Testing,
a general testing paradigm enabling repository-level correctness evaluation without the requirement of formal specifications.
A testing agent operates within a sandboxed environment on top of its own copy of the repository, and then translates informal properties into analysis and evidence gathering actions, finally returning a \textsc{pass}/\textsc{fail}/\textsc{inconclusive} decision along with an evidence report.
We evaluate agentic testing on three bug domains: ML pipeline bugs, security vulnerabilities, and citation hallucinations in \LaTeX{} documents, comparing against domain-specific analyses and general-purpose code reviewing agents.
\end{abstract}

\section{Introduction}
% Outline:
% 1. Media level excitement.
% -  Software correctness is a critical problem.
%
% 2. New development. Research level excitement.
% - Software engineering is being reshaped by highly capable LLM agents. They
%   are able to understand complex codebases and make changes based on high level descriptions.
%
% 3. Problem
% - Even for human written code, it is very challenging to ensure its correctness.
% - Finding bugs requires extensive skill and effort.
%
% 4. Insufficiency of existing methods
% - Program analysis methods exist, but they require formal specs of the bugs making them hard to adapt to new codebases and bugs.
%
% 5. Big picture problem
% - How do we evaluate the correctness of programs in regards to high level rather than formal specs?
%
% 6. Our approach
% - We propose a general testing agent for evaluating repository-level code correctness without requiring formal specs.
% - We perform an evaluation on three bug domains: ML pipeline bugs, security vulnerabilities, and citation hallucinations in latex documents.

% Main figure: I'm thinking we can highlight an interesting bug which the agent finds that would clearly illustrate the capabilities of the method.
% I think one of the ML pipeline bugs from an ICLR submission would stand out.

Repository-scale software systems combine code, data, configuration, and documentation, and their correctness is often specified only in natural language either explicitly or implicitly based on the intentions of the developers.
Validating such intent remains difficult and costly in practice \citep{zeller2009why}.
Even in secure-development contests where teams are explicitly incentivized to build correct, secure systems, submissions routinely contain vulnerabilities \citep{ruef2016bibifi}.
% In this work we focus on repositories where high-level intent must be evaluated without formal specifications.

At the same time, software engineering is being reshaped by highly capable AI agents.
Benchmarks such as SWE-bench~\citep{jimenez2024swebench} require models to resolve real-world bug reports in repositories, and recent analyses show that agentic workflows are becoming increasingly effective at solving these repository-level tasks~\citep{jimenez2024swebench,xia2024agentless,guo2025repoaudit}.

Existing approaches to correctness provide complementary but incomplete coverage.
Static program analysis can be sound with respect to a formal semantics, but it typically requires explicit specifications and still trades completeness for tractability \citep{cousot1979systematic, nielson1999principles}.
Recent LLM-assisted query systems such as IRIS~\citep{li2024iris}, QLCoder~\citep{wang2025qlcoder}, and KNighter~\citep{yang2025knighter} reduce manual effort but still depend on formalized properties which can fail to generalize and miss higher level issues.
Testing and fuzzing explore concrete executions and can uncover defects, but they again require extensive manual effort to extend to new properties or repositories for validating~\citep{manes2018fuzzing}.
Together, these methods leave a gap for evaluating high-level intent when specifications are informal.

This gap raises a broader question.
How do we evaluate program correctness with respect to high-level intent rather than formal specifications?
Answering this question requires methods that translate informal intent into actionable, repository-level testing procedures.

To address this gap, we introduce agentic testing for repository-scale correctness.
The agent takes a repository and a natural language property specification, performs analysis and evidence gathering, and outputs a \textsc{pass}/\textsc{fail}/\textsc{inconclusive} decision with an evidence report.
The design emphasizes repository-level context and allows the agent to execute and modify code within a sandboxed copy of the repository, without persisting changes.

We evaluate agentic testing on three bug domains: ML pipeline bugs, security vulnerabilities, and citation hallucinations in \LaTeX{} documents.
This evaluation spans multiple datasets and compares the agent against domain-specific analyses and general-purpose code review agents.

\begin{figure*}[t]
  \centering
  \includegraphics[width=0.95\linewidth]{figures/agentic_testing_evidence_verification.png}
  \caption{Agentic testing evaluates a repository against an informal property by gathering evidence, issuing a \textsc{pass}/\textsc{fail}/\textsc{inconclusive} decision, and exposing the evidence for downstream audit.}
  \label{fig:motivating-bug}
\end{figure*}

We summarize our contributions as follows.
\begin{itemize}
  \item We formulate repository-level correctness as testing informal properties specified in natural language, motivating a setting where high-level intent must be evaluated without formal specifications.
  \item We propose an agentic testing system that operates over repositories and produces \textsc{pass}/\textsc{fail}/\textsc{inconclusive} decisions with evidence reports.
  \item We construct evaluation datasets across three domains---ML pipelines, security vulnerabilities, and citation hallucinations---and evaluate the agent against domain-specific baselines and model backbones using both accuracy and evidence quality.
\end{itemize}



\section{Informal Properties of Code}
% In this section, we introduce the three domains that we evaluate agentic testing on, and how we constructed the datasets.
% Datasets:
% 1. ML pipeline bugs
% - Dataset 1: We scrape 50 random submissions from each of three Kaggle competitions: Titanic (tabular), Diabetic Retinopathy (image), and NLP with Disaster Tweets (text).
% - Dataset 2: We use 100 randomly selected submissions to ICLR 2026 that include supplementary material containing at least 1 python file. We additionally include the paper pdf in the repository.
% 2. Security vulnerabilities
% - Dataset 1: We use the Fall 2025 competition from the build-it-break-it-fix-it (BIBIFI) contest. This contains 25 submissions of programs implementing a secure ATM along with annotations of verified vulnerabilities.
% - Dataset 2: We use the CWE-Bench dataset from the IRIS paper.
% 3. Citation hallucinations in latex documents
% - Dataset 1: We scrape all ICLR 2026 submissions to OpenReview that were also available on arXiv (with the same title) and download their complete latex source repository.

% Natural language properties are specified in the following markdown files:
% 1. ML pipeline bugs (properties/ml-properties.ml)
% 2. Security vulnerabilities (properties/bibifi-properties.ml and properties/cwe-properties.ml)
% 3. Citation hallucinations in latex documents (properties/hallucination-properties.ml)
Software correctness is often judged by informal properties that are stated at a high level rather than encoded as formal specifications.
Static analysis tools require explicit, formalized queries, and recent LLM-assisted methods aim to reduce the burden of writing those queries, but the target properties remain informal and domain-specific \citep{codeql,li2024iris,wang2025qlcoder}.
We define an informal property as a natural language statement about expected behavior or artifacts in a repository.
Such properties are central in ML pipelines, security auditing, and scientific writing, where failures are often detectable only through repository-level context and human intent \citep{schoop2021umlaut,ruef2016bibifi,ye2023agree,haan2025semanticcite}.

This section motivates the problem setting and introduces the datasets used in our study.
We focus on three domains of informal properties: ML pipeline bugs, security vulnerabilities, and citation hallucinations in \LaTeX{} documents.
We provide representative examples below and include the full property lists in \Cref{app:properties}.

We formalize the setting as follows.
Classical program analysis asks whether a formal property holds on all traces of a program.
Let $T(R)$ denote the set of concrete execution traces induced by repository $R$, and let $\phi$ be a formal predicate over traces, such as ``no array out of bounds.''
The traditional goal is to decide whether $\forall t \in T(R), \phi(t)$ holds, or to find a counterexample trace when it does not.
In our setting, the key difference is that the property is specified informally rather than as a formal predicate.
Let $p$ be a natural language statement about expected behavior or artifacts in $R$, such as ``no training on the test set'' or ``every cited DOI resolves.''
Agentic testing evaluates a pair $(R, p)$ and returns a decision $y \in \{\textsc{pass}, \textsc{fail}, \textsc{inconclusive}\}$ with an evidence report $E$.
The evidence report should identify the relevant files, executions, and outputs that justify $y$ and enable independent verification.
This framing treats properties as inputs rather than handcrafted queries, enabling evaluation of high-level intent without requiring formal specifications.

\subsection{ML Pipeline Bugs}
ML pipeline bugs are often subtle and are diagnosed using domain heuristics and end-to-end context \citep{schoop2021umlaut}.
We use two datasets of ML pipeline repositories.
The first dataset contains 50 random submissions from each of three Kaggle competitions.
The competitions are Titanic (tabular), Diabetic Retinopathy (image), and NLP with Disaster Tweets (text).
The second dataset contains 100 randomly selected ICLR 2026 submissions with supplementary material that includes at least one Python file.
For those submissions, we include the paper PDF in the repository.
Representative ML properties include:
\begin{itemize}
  \item No leakage from test to train/val sets.
  \item Data augmentation is applied only to training data.
  \item Reported metrics use the correct split and definitions.
\end{itemize}

\subsection{Security Vulnerabilities}
Security vulnerabilities are a canonical setting where correctness is defined by informal threat models and high-level intent \citep{ruef2016bibifi,li2024iris}.
We use two datasets of security-focused code.
The first dataset is the Fall 2025 build-it-break-it-fix-it competition.
It contains 25 submissions of programs implementing a secure ATM with annotations of verified vulnerabilities.
The second dataset is the CWE-Bench dataset from the IRIS paper.
Representative security properties include:
\begin{itemize}
  \item Authenticated integrity for all messages.
  \item Replay protection is correct.
  \item Inputs are bounds-checked before processing.
\end{itemize}

\subsection{Citation Hallucinations in \LaTeX{} Documents}
Citation accuracy is another informal property where the intended claim-to-source relationship is rarely formalized \citep{ye2023agree,haan2025semanticcite}.
We use a dataset of \LaTeX{} repositories derived from ICLR 2026 submissions.
We scrape all ICLR 2026 submissions to OpenReview that were also available on arXiv with the same title.
We download the complete \LaTeX{} source repository for each paper.
Representative citation properties include:
\begin{itemize}
  \item Every listed DOI must resolve to the cited paper.
  \item Every cited arXiv identifier must correspond to the claimed paper.
  \item Every URL is accessible or archived in the Wayback Machine.
\end{itemize}

\begin{table*}[t]
  \centering
  \caption{Dataset summary for the current benchmark package. Kaggle counts are real repositories; synthetic Kaggle contains injected repositories derived from real Kaggle submissions.}
  \label{tab:datasets}
  \begin{tabular}{llll}
    \toprule
    Domain & Dataset & Properties & Size \\
    \midrule
    ML pipeline bugs & Synthetic Kaggle & 15 & 75 repos / 1125 pairs \\
    ML pipeline bugs & Real Kaggle & 16 & 120 repos / 1920 pairs \\
    ML pipeline bugs & ICLR '26 & -- & Planned \\
    Security vulnerabilities & BIBIFI Fall '25 & 8 & 25 submissions \\
    Security vulnerabilities & CWE-Bench & 4 & 197 cases \\
    Citation hallucinations & ICLR '26 arXiv & -- & Planned \\
    \bottomrule
  \end{tabular}
\end{table*}

\section{Agentic Testing}
% In this section, we describe our proposed agentic testing system.
We propose a testing agent that evaluates whether an informal property holds for a given repository.
The agent takes two inputs: the code repository and a property specification written in natural language.
It produces a decision of \textsc{pass}, \textsc{fail}, or \textsc{inconclusive}, along with an evidence report that justifies the decision.
The design emphasizes repository-level reasoning rather than local code snippets and does not require formal specifications.

\begin{figure}[t]
  \centering
  \includegraphics[width=\linewidth]{figures/agentic_testing_evidence_verification.png}
  \caption{Overview of agentic testing.
  The agent ingests a repository and a property specification, executes three phases of analysis, and outputs a decision with supporting evidence.}
  \label{fig:overview}
\end{figure}

\paragraph{Phase 1: Analysis.}
The agent first builds a high-level model of the repository to situate the property.
It identifies key files, entry points, and relevant artifacts such as configuration files, datasets, and documentation.
When applicable, it assesses a domain-specific threat model or assumptions that shape the property (e.g., attacker capabilities in security settings).
This phase narrows the search space and yields a plan for evidence gathering.

\paragraph{Phase 2: Evidence gathering.}
The agent reviews code, tests, and data processing logic to find evidence for or against the property.
It can run code when necessary to validate behavior, reproduce results, or check for violations.
The goal is to collect concrete, reproducible evidence tied to specific files, logs, or outputs.

\paragraph{Phase 3: Decision.}
The agent aggregates the gathered evidence and produces a decision.
\textsc{Pass} indicates the property is supported by evidence, \textsc{fail} indicates a violation, and \textsc{inconclusive} indicates insufficient evidence.
The agent outputs an evidence report that highlights the relevant files, outputs, and reasoning steps that support the classification.
The agent operates in a sandboxed copy of the repository and may modify code to probe properties, but these modifications do not persist to the original repository.

\section{Experiments}
\subsection{Setup}

% Baselines:
% - Vulnerabilities: CodeQL, BugBog
% - Citation hallucinations: Existing approach for finding hallucinated citations, BugBot
% - ML pipeline bugs: UMLAUT, BugBot

% Models:
% - GPT-5-mini
% - Claude-4.5-Sonnet
% - GPT-5.2-codex

% Evaluation:
% - We evaluate both classification (buggy or not) of the (repo, property) pairs
%   as well as evaluate the evidence provided by the methods to determine if their
%   classifications are properly justified.
% - For datasets with ground truth (CWE-Bench and citation hallucination), we
%   compute both recall and precision. For datasets without ground truth, we only
%   compute the precision using human evaluation. Where possible, each property
%   is evaluated by three independent evaluators to ensure high quality. This evaluation
%   takes 2-3 minutes per (repo, property) pair.
We evaluate baselines for each domain using the following domain-specific methods.
For ML bugs, we compare with TrainCheck \citep{traincheck}.
For vulnerabilities, we include CodeQL \citep{codeql}.
For citation hallucinations, we include a prior detection approach for hallucinated references \citep{agrawal2023hallucinated}.
We also compare with BugBot \citep{cursor_bugbot} and Codex Code Review \citep{openai_codex} as general-purpose, open-ended code review baselines.
We refer to these as Code Review (CR) baselines because they are prompted to find any bugs, without conditioning on a specific property.
To indicate the specific model family, we denote Codex Code Review with a model suffix, e.g., CR-gpt-5.2-codex, while BugBot remains CR (BugBot).
Our method, Agentic Testing (AT), instead conditions on a property and focuses evidence gathering on that target.

We evaluate the testing agent with GPT-5-mini, Claude-4.5-Sonnet, and GPT-5.2-codex \citep{openai_models,anthropic_models}.
For evaluation, we collect the classification and evidence provided by each method on each (repo, property) pair and then either use ground truth or human evaluation to determine the correctness of any bug detections.
This human evaluation takes 2--3 minutes per (repo, property) pair.
This enables us to assess the precision of each method. For datasets with ground truth (CWE-Bench and citation hallucination), we also compute recall.

\subsection{RQ1: Classification Accuracy vs. Domain Baselines}
We ask whether the agent can accurately classify (repo, property) pairs compared to domain baselines.
For datasets with ground truth, we report precision and recall for each method.
For datasets without ground truth, we report precision under human evaluation.
We present results by domain and dataset to expose where the agent is strongest or weakest.
For security vulnerabilities, where ground-truth vulnerability labels are available, \Cref{tab:rq1} reports verified recall, abstention rate, and cost efficiency.
For domains without complete ground truth, \Cref{tab:rq1} reports verified precision from human validation.
AT achieves the strongest verified recall on both vulnerability subdatasets.

\begin{table*}[t]
  \centering
  \caption{Current headline results. For Synthetic Kaggle, coverage is the non-\textsc{inconclusive} rate and macro F1 is measured against injected labels. For Real Kaggle, macro F1 is estimated from the conservative 15\% fail audit. Security rows retain verified recall values from the current draft.}
  \label{tab:rq1}
  \begin{tabular}{llrrcc}
    \toprule
    Subdataset & Method & Coverage & Macro F1 / Recall & Abstention Rate & Notes \\
    \midrule
    \multicolumn{6}{c}{\textbf{ML pipeline bugs}} \\
    \multirow{4}{*}{Synthetic Kaggle}
      & TrainCheck & 0.000 & -- & 1.000 & attempted baseline \\
      & Reviewer mode 0 & 0.264 & 0.669 & 0.736 & Qwen reviewer \\
      & Reviewer mode 1 & 0.321 & 0.722 & 0.679 & Qwen reviewer \\
      & AT-Qwen, 10 examples & 0.903 & 0.813 & 0.097 & best coverage \\
    \midrule
    \multirow{4}{*}{Real Kaggle}
      & Direct property, Qwen Flash & 0.611 & -- & 0.389 & no full labels \\
      & Snapshot reviewer, GPT-5-mini & 0.111 & -- & 0.889 & conservative CR baseline \\
      & AT-Qwen, 10 examples & 0.869 & 0.933 & 0.131 & conservative audit \\
      & AT-Qwen, 20 examples & 0.871 & 0.938 & 0.129 & conservative audit \\
    \midrule
    \multicolumn{6}{c}{\textbf{Security vulnerabilities}} \\
    \multirow{5}{*}{BIBIFI}
      & CodeQL & -- & 1.75\% & 0.00\% & N/A \\
      & Codex-gpt-5-mini & -- & 29.82\% & 70.00\% & 0.003 \\
      & Codex-gpt-5.2 & -- & -- & -- & -- \\
      & AT-gpt-5-mini & -- & 64.91\% & 0.00\% & 0.014 \\
      & AT-gpt-5.2 & -- & \textbf{71.93\%} & 0.00\% & 0.119 \\
      \midrule
    \multirow{4}{*}{CWE-Bench}
      & CodeQL & -- & 16.33\% & 0.00\% & N/A \\
      & Codex-gpt-5-mini & -- & 6.25\% & 69.39\% & 0.011 \\
      & AT-gpt-5-mini & -- & 6.12\% & 15.82\% & 0.052 \\
      & AT-gpt-5.2 & -- & \textbf{24.49\%} & 23.47\% & 0.061 \\
    \bottomrule
  \end{tabular}
\end{table*}

\subsection{RQ2: Evidence Quality}
We ask how strong the evidence is that the agent provides for its classifications relative to baselines.
We use the same human evaluation protocol as in the setup to assess evidence quality.
We report evidence quality alongside accuracy to distinguish correct predictions from well-justified predictions.
For ML pipeline bugs, the audit separates prediction quality from evidence quality.
On real Kaggle, the conservative fail audit verifies 122 of 152 sampled predicted failures as true failures.
On synthetic Kaggle, a disagreement audit finds that 29 of 53 usable sampled disagreements are model misses, while 24 reflect ground-truth or property-definition mismatches.
This suggests that the synthetic benchmark is useful but pessimistic for natural-language properties whose boundaries are partly interpretive.

\begin{table}[t]
  \centering
  \caption{Evidence audit summary for the ML pipeline experiments.}
  \label{tab:rq2}
  \begin{tabular}{lll}
    \toprule
    Audit & Result & Notes \\
    \midrule
    Real Kaggle fail audit & 122/152 true failures & conservative 15\% sample \\
    Synthetic disagreement audit & 29/53 model misses & usable sampled disagreements \\
    Synthetic label audit & 24/53 definition mismatches & property-boundary ambiguity \\
    \bottomrule
  \end{tabular}
\end{table}

\subsection{RQ3: Variation Across Domains and Models}
We ask how performance varies across domains and model backbones.
We compare results for GPT-5-mini, Claude-4.5-Sonnet, and GPT-5.2-codex across the three domains.
This analysis isolates whether gains are driven by domain structure or by model capability.
The current ML pipeline package uses Qwen 35B as the main AT backbone, Qwen Flash for the direct-property baseline, and GPT-5-mini for the snapshot reviewer baseline.
We keep this model mismatch explicit: the direct-property baseline is a useful auxiliary baseline, but it is not a strict model-matched comparison to AT-Qwen.

\begin{table*}[t]
  \centering
  \caption{Model/domain status for current runs.}
  \label{tab:rq3}
  \begin{tabular}{llll}
    \toprule
    Model & Domain & Status & Main use \\
    \midrule
    \multicolumn{4}{l}{\textbf{ML pipeline bugs}} \\
    Qwen 35B & ML pipeline bugs & complete & AT main results \\
    Qwen Flash & ML pipeline bugs & complete & direct-property baseline \\
    GPT-5-mini & ML pipeline bugs & complete & snapshot reviewer baseline \\
    \addlinespace
    \multicolumn{4}{l}{\textbf{Security vulnerabilities}} \\
    GPT-5-mini & Security vulnerabilities & partial & AT and CR rows \\
    GPT-5.2 & Security vulnerabilities & partial & stronger AT row \\
    GPT-5.2-codex & Security vulnerabilities & partial & code-review row \\
    \addlinespace
    \multicolumn{4}{l}{\textbf{Citation hallucinations}} \\
    GPT-5-mini & Citation hallucinations & planned & AT and CR rows \\
    Claude-4.5-Sonnet & Citation hallucinations & planned & robustness \\
    GPT-5.2-codex & Citation hallucinations & planned & code-review row \\
    \bottomrule
  \end{tabular}
\end{table*}

\section{Related Work}
% Program analysis
% - CodeQL, ML pipeline testing methods
%
% Agents for evaluation
% - IRIS, RepoAudit, BugScope, BugBot (cursor), Codex Code Review
\paragraph{Static analysis and specification burden.}
Query-based static analysis systems such as CodeQL encode bug patterns as declarative queries and support custom project-specific checks, but they require specialized expertise to author and maintain \citep{codeql}.
Recent work uses LLMs to reduce this burden, including QLCoder for synthesizing CodeQL queries from vulnerability metadata \citep{wang2025qlcoder}.
IRIS combines LLMs with static analysis to infer taint specifications and perform whole-repository vulnerability detection without relying solely on human-labeled specifications \citep{li2024iris}.
KNighter generates static analysis queries from examples, further illustrating the shift from hand-written analysis rules toward learned or generated checks \citep{yang2025knighter}.

\paragraph{Repository-level agents and evaluation.}
RepoAudit proposes an autonomous LLM agent for repository-level auditing that reasons over data-flow facts and validates candidate reports \citep{guo2025repoaudit}.
BugScope learns bug patterns from examples and applies multi-agent auditing to detect diverse defects \citep{liu2025bugscope}.
SWE-bench frames repository-scale issue resolution and highlights the difficulty of end-to-end fixes that require coordinated edits across files \citep{jimenez2024swebench}.
Agentless shows that simplified localization--repair pipelines can be competitive on SWE-bench Lite, providing a strong baseline for agentic approaches \citep{xia2024agentless}.

\paragraph{Domain-specific debugging and citation grounding.}
ML debugging tools such as Umlaut use expert heuristics to surface errors in deep learning programs by linking code structure and model behavior \citep{schoop2021umlaut}.
For citation quality, AGREE targets grounded citation generation and SemanticCite focuses on citation verification with evidence-based reasoning \citep{ye2023agree,haan2025semanticcite}.
These domain-specific efforts motivate a general approach to informal properties that spans code and documents \citep{schoop2021umlaut,ye2023agree,haan2025semanticcite}.


\section{Conclusion}
Agentic testing reframes repository-level correctness evaluation as checking informal natural-language properties against code, data, and artifacts.
The ML pipeline experiments show that property-conditioned agents can produce substantially higher usable coverage than generic review-style baselines while still exposing evidence for audit.
At the same time, the synthetic audit highlights an important limitation: natural-language properties are not always perceived identically by dataset construction, model predictions, and human verification.
This makes evidence-centered auditing part of the benchmark rather than a post-hoc detail.
Future work should expand model-matched baselines, complete the security and citation domains, and refine protocols for adjudicating property-definition disagreements.

\section*{Impact Statement}
This work aims to improve the reliability of software and machine-learning repositories by making informal correctness properties easier to test and audit.
The same capability could be used to discover security flaws or implementation weaknesses in real systems.
We therefore emphasize sandboxed execution, evidence reporting, and responsible disclosure norms for any deployment on non-consenting repositories.

% In the unusual situation where you want a paper to appear in the
% references without citing it in the main text, use \nocite
% \nocite{langley00}

\bibliography{vibetest}
\bibliographystyle{icml2026}

%%%%%%%%%%%%%%%%%%%%%%%%%%%%%%%%%%%%%%%%%%%%%%%%%%%%%%%%%%%%%%%%%%%%%%%%%%%%%%%
%%%%%%%%%%%%%%%%%%%%%%%%%%%%%%%%%%%%%%%%%%%%%%%%%%%%%%%%%%%%%%%%%%%%%%%%%%%%%%%
% APPENDIX
%%%%%%%%%%%%%%%%%%%%%%%%%%%%%%%%%%%%%%%%%%%%%%%%%%%%%%%%%%%%%%%%%%%%%%%%%%%%%%%
%%%%%%%%%%%%%%%%%%%%%%%%%%%%%%%%%%%%%%%%%%%%%%%%%%%%%%%%%%%%%%%%%%%%%%%%%%%%%%%
\newpage
\appendix
\onecolumn
\section{Full Informal Property Lists}
\label{app:properties}
\subsection{ML Pipeline Properties}
\begin{itemize}
  \item No leakage from test to train/val sets.
  \item If there is any model selection or hyperparameter tuning, then it uses val only.
  \item If there is any model training, then all model parameters are updated during training (no frozen layers unless explicitly intended).
  \item Data is loaded and preprocessed correctly (e.g., no all-black images, no text with weird or unexpected characters, tables look correct and feature values are reasonable).
  \item Any data augmentation is only performed on the training dataset; val/test dataset use deterministic preprocessing.
  \item Any class-imbalance handling (weighted loss / sampler) applies only to training data.
  \item No path/filename text is fed as a feature unless explicitly intended.
  \item Reported metrics use the correct split(s) and the exact definitions claimed (e.g., micro vs macro, top-k).
  \item Randomizing the labels results in accuracy dropping to be near a random guessing baseline (may not be 0.5 if the data is imbalanced) on a validation set.
  \item The model after the full training procedure outperforms a simple baseline (e.g., random or majority class) on an evaluation set.
  \item The model can overfit a single (or tiny) batch to near-zero loss.
  \item Training loss generally decreases during training and plateaus within the number of epochs used (if no loss is logged, then add logging to check this).
  \item Training dataloader shuffles or training dataset is shuffled for training; val/test do not shuffle.
  \item Accuracy should be deterministic (same value) when running the model evaluation multiple times without retraining.
  \item Model and inputs are consistently moved to one device; no implicit CPU<->GPU transfers; dtype policy (e.g., fp32/bf16) is applied consistently.
  \item Code does not include any secrets (API keys, passwords, etc.).
  \item Code does not use explicit loops or Python control flow when matrix operations could have been used to implement the exact same operation more efficiently.
\end{itemize}

\subsection{Security Properties}
\begin{itemize}
  \item Complete transport confidentiality: All sensitive protocol payload data is encrypted end-to-end on every request/response, with no plaintext fields that depend on secrets.
  \item No sensitive leakage via metadata/side channels: Observable metadata (e.g., connection info, timing/length patterns) does not vary with secrets or internal outcomes beyond an explicitly allowed constant envelope.
  \item Authenticated integrity for all messages: Every message is protected with a keyed integrity mechanism (e.g., AEAD or MAC) and verification is enforced for every message type, including error responses.
  \item Security checks are enforced: All required validations (including parsing/extension checks) are performed and any failure causes the message/session to be rejected.
  \item Replay protection is correct: Each message uses freshness (nonce/timestamp) that is verified, tracked (no reuse), and rejected when stale, duplicated, or out of policy.
  \item IVs are unique and unpredictable: Every encryption operation uses an IV/nonce that is never reused and is not predictable from public information.
  \item Secrets are high-entropy and non-derivable: Security-critical values are generated with sufficient randomness or derived from secret keys, and are not deterministically computable from public identifiers.
  \item Client authentication is required: Clients must authenticate before any privileged interaction, and unauthenticated clients cannot successfully perform protocol actions.
  \item Inputs are bounds-checked before processing: All incoming messages are length/bounds validated against protocol limits before parsing or allocation.
  \item Robust error handling and state consistency: Exceptions are handled safely and any desynchronization between client/server state triggers a safe reset/reject rather than continuing in an inconsistent state.
\end{itemize}

\subsection{Citation Properties}
\begin{itemize}
  \item Every listed arxiv id of the form yymm.nnnnn associated with a reference corresponds to an accessible arxiv webpage at http://arxiv.org/abs/\texttt{<identifier>} for the paper with the same title and authors as cited.
  \item Every listed DOI (Digital Object Identifier) must resolve via https://doi.org/\texttt{<doi>} to a landing page for the specific paper cited, matching the title and authors listed.
  \item If a paper is cited as being published in ``Conference X, Year Y,'' the official proceedings of Conference X for Year Y must contain that paper.
  \item Every url can be accessed without a 404 response or exists in the Wayback Machine.
\end{itemize}
%%%%%%%%%%%%%%%%%%%%%%%%%%%%%%%%%%%%%%%%%%%%%%%%%%%%%%%%%%%%%%%%%%%%%%%%%%%%%%%
%%%%%%%%%%%%%%%%%%%%%%%%%%%%%%%%%%%%%%%%%%%%%%%%%%%%%%%%%%%%%%%%%%%%%%%%%%%%%%%

\end{document}

% This document was modified from the file originally made available by
% Pat Langley and Andrea Danyluk for ICML-2K. This version was created
% by Iain Murray in 2018, and modified by Alexandre Bouchard in
% 2019 and 2021 and by Csaba Szepesvari, Gang Niu and Sivan Sabato in 2022.
% Modified again in 2023 and 2024 by Sivan Sabato and Jonathan Scarlett.
% Previous contributors include Dan Roy, Lise Getoor and Tobias
% Scheffer, which was slightly modified from the 2010 version by
% Thorsten Joachims & Johannes Fuernkranz, slightly modified from the
% 2009 version by Kiri Wagstaff and Sam Roweis's 2008 version, which is
% slightly modified from Prasad Tadepalli's 2007 version which is a
% lightly changed version of the previous year's version by Andrew
% Moore, which was in turn edited from those of Kristian Kersting and
% Codrina Lauth. Alex Smola contributed to the algorithmic style files.
